
\subsubsection{5.1 Main Results}

\begin{table}[htbp]
\centering
\resizebox{\columnwidth}{!}{%
\begin{tabular}{llll}
\toprule
Method & AUROC & Threshold & Status \\
\midrule
Random baseline & 0.50 & --- & Reference \\
Choice Entropy & 0.7703 & 0.55 & PASS \\
Max Probability & 0.8068 & 0.55 & PASS \\
Semantic Entropy & 0.5645 & 0.70 & FAIL \\
\bottomrule
\end{tabular}
}
\end{table}

\textbf{Finding 1:} Token-level methods discriminate hallucinations. Max probability (AUROC = 0.8068) and choice entropy (AUROC = 0.7703) both exceed the 0.55 threshold.

\textbf{Finding 2:} Semantic entropy underperforms. AUROC = 0.5645, below the 0.70 threshold and 0.24 points below max probability.

\subsubsection{5.2 Mechanism Verification}

\begin{table}[htbp]
\centering
\resizebox{\columnwidth}{!}{%
\begin{tabular}{llll}
\toprule
Check & Value & Expected & Status \\
\midrule
Avg clusters per question & 4.92 & < 5.0 & PASS \\
Entropy standard deviation & 0.0752 & > 0 & PASS \\
\bottomrule
\end{tabular}
}
\end{table}

\textbf{Finding 3:} The semantic clustering mechanism functions correctly. Clustering occurs (average 4.92 clusters from 5 samples) and entropy varies across questions.

\subsubsection{5.3 Format-Dependency Analysis}

The resolution to the apparent contradiction---correct mechanism but poor discrimination---lies in the task format. Multiple-choice responses are single letters (''A'', ''B'', ''C'', ''D''). NLI models cannot meaningfully compare ''A'' versus ''B'' for entailment. Each response tends to form its own cluster (average 4.92 from 5 samples). Entropy over near-singleton clusters provides no discriminative signal.

Token-level methods extract uncertainty directly from logits. They do not require semantic comparison and function regardless of answer length.

